\documentclass[11pt,a4paper]{article}
\usepackage[utf8]{inputenc}
\usepackage[T1]{fontenc}
\usepackage[spanish,es-nodecimaldot]{babel}
\usepackage[margin=2.3cm]{geometry}
\usepackage{amsmath,amssymb,booktabs,array,tabularx}
\usepackage{xcolor}
\usepackage{enumitem}
\usepackage{tikz}
% Compatibilidad con los caracteres activos de babel en español.
\usetikzlibrary{babel,arrows.meta,positioning}
\usepackage{listings}
\usepackage{hyperref}
\hypersetup{colorlinks=true,linkcolor=blue!45!black,urlcolor=blue!45!black}
\setlength{\parindent}{0pt}
\setlength{\parskip}{5pt}
\setlist{itemsep=3pt,topsep=4pt}
\renewcommand{\arraystretch}{1.18}
\newcommand{\code}[1]{\texttt{\detokenize{#1}}}
\newcommand{\casilla}{$\square$\ }
\lstset{basicstyle=\ttfamily\small,breaklines=true,breakatwhitespace=true,frame=single,columns=fullflexible,keepspaces=true}
\emergencystretch=3em
\title{Puesta en marcha de los encoders magnéticos AS5600\newline (fuente de posición 3)}
\author{Instrucciones para cuando se disponga de los sensores\\Proyecto MiYoke PD V2.6}
\date{22 de octubre de 2026}
\begin{document}
\maketitle

\begin{center}
\fbox{\parbox{0.9\textwidth}{\textbf{Idea principal.} El firmware V2.6 ya incluye la lectura de dos AS5600 por I2C (\code{YOKE_POSITION_SOURCE=3}), probada solo con pruebas unitarias. Este documento ordena lo que falta hacer con los sensores en la mano: probarlos en banco, montarlos con la distancia correcta al imán, medir el centro y los extremos, cargar esos valores y ensayar fallas y tiempos antes de usar los motores.}}
\end{center}

\textbf{Alcance.} Procedimiento a ejecutar por el usuario. Los valores numéricos de \code{mis_defines.h} indicados como estimación (cuentas por extremo de pitch) deben reemplazarse por los medidos. La referencia técnica completa está en \code{LEEME md/LEEME_AS5600.md} y el diseño en \code{LEEME md/DISENO_encoder_as5600.md}.

\tableofcontents

\section{Cómo funciona el sensor (lo necesario para entender los pasos)}
\begin{itemize}
\item El AS5600 mide el ángulo del campo de un imán diametral que gira sobre el chip. Entrega 4096 cuentas por vuelta (\(360/4096 \approx 0{,}088^\circ\) por cuenta) y es \textbf{absoluto dentro de una vuelta}: no cuenta vueltas. Si el eje gira más, las vueltas las acumula el firmware.
\item La dirección I2C es fija (0x36). Por eso roll usa un bus propio (I2C1) y pitch comparte el bus del LCD (I2C0).
\item El chip regula su ganancia (AGC) según la intensidad del campo. La telemetría la muestra como \code{g=}; el valor óptimo está en el centro del rango, cerca de 64 a 3,3~V. Esa es la herramienta para ajustar la distancia del imán.
\item La telemetría serie (115200 baudios, una línea cada 100~ms alternando Roll/Pitch) tiene la forma
\begin{lstlisting}
P a=2731 s=OK g=63 x=1250 v=0 f=0 p=0 | R a=1022 s=OK g=66 ...
\end{lstlisting}
\code{a} es el ángulo crudo (0..4095), \code{s} el estado, \code{g} el AGC, \code{x} la posición \(\times 10000\), \code{v} la velocidad, \code{f} el comando y \code{p} el PWM.
\end{itemize}

\begin{center}
\begin{tabularx}{\textwidth}{lX}
\toprule
Estado \code{s=} & Significado y acción\\
\midrule
\code{OK} & Lectura correcta.\\
\code{DEB} & Campo débil: imán demasiado lejos o desalineado. Acercar.\\
\code{FUE} & Campo fuerte: imán demasiado cerca. Alejar.\\
\code{SIN} & Imán no detectado. Revisar montaje; si el módulo nunca informa el bit MD con el imán bien puesto, ver sección~\ref{sec:problemas}.\\
\code{I2C} & El sensor no responde: cableado, alimentación o dirección.\\
\code{BLQ} & Bloqueo por posible vuelta perdida: fuerza en cero hasta reiniciar la Pico.\\
\bottomrule
\end{tabularx}
\end{center}

\section{Materiales y verificación previa}
\begin{itemize}
\item[\casilla] Dos módulos AS5600 con su imán diametral (normalmente 6~mm de diámetro por 2,5 a 3~mm).
\item[\casilla] Soportes del sensor de aluminio, plástico o impresión 3D; adaptadores del imán al eje de material no magnético.
\item[\casilla] Cable de 4 hilos corto, idealmente trenzado (SDA con GND, SCL con GND).
\item[\casilla] Galgas o separadores de 0,5, 1, 1,5 y 2~mm para fijar la distancia imán--chip.
\item[\casilla] Confirmar a qué tensión están las pull-ups del bus del LCD. Si el módulo PCF8574 se alimenta a 5~V con pull-ups a 5~V, el AS5600 de pitch se conecta \textbf{del lado de 3,3~V} del adaptador de nivel (ver \code{LEEME_PICO2.md}). El AS5600 a 3,3~V no tolera 5~V en SDA/SCL.
\item[\casilla] Si el módulo tiene puente o selector de alimentación 3,3/5~V, dejarlo en 3,3~V.
\end{itemize}

\section{Conexión eléctrica}
\begin{center}
\begin{tabular}{lll}
\toprule
Pin del módulo & AS5600 de roll & AS5600 de pitch\\
\midrule
VCC & 3V3(OUT), pin físico 36 & 3V3(OUT), pin físico 36\\
GND & GND & GND\\
SDA & GP14, pin físico 19 & GP4, pin físico 6 (junto al LCD)\\
SCL & GP15, pin físico 20 & GP5, pin físico 7 (junto al LCD)\\
DIR & GND & GND\\
OUT, PGO/GPO & sin conectar & sin conectar\\
\bottomrule
\end{tabular}
\end{center}
\begin{itemize}
\item \textbf{DIR no debe quedar suelto}: si flota, el sentido de giro informado puede cambiar al azar.
\item Mantener los cables I2C lejos de los cables del motor y del puente H; no tenderlos juntos en el mismo mazo.
\item No subir el bus del LCD a 400~kHz: el PCF8574 está especificado para 100~kHz.
\end{itemize}

\section{Paso 1: prueba en banco, sin montar}
Objetivo: comprobar cada módulo antes de pensar en la mecánica. \textbf{Motores sin alimentación} (fuente de los BTS7960 apagada).
\begin{enumerate}
\item Compilar la variante AS5600:
\begin{lstlisting}
powershell -ExecutionPolicy Bypass -File herramientas/compilacion/probar.ps1 -As5600
\end{lstlisting}
El UF2 queda en \code{build-pico2-as5600/}. Cargarlo en la Pico con BOOTSEL.
\item Conectar ambos sensores (si se conecta uno solo, el otro eje mostrará \code{s=I2C} y la fuerza quedará en cero: es lo esperado).
\item Abrir el monitor serie a 115200 baudios. Con el simulador cerrado se ven las líneas de movimiento.
\item Sostener el imán con la mano a 1--2~mm sobre el chip, centrado, y girarlo lentamente:
  \begin{itemize}
  \item[\casilla] \code{s=OK} con el imán puesto y \code{s=SIN} (o \code{DEB}) al retirarlo.
  \item[\casilla] \code{a=} recorre 0..4095 suavemente en una vuelta, sin saltos, y pasa de 4095 a 0 al completarla.
  \item[\casilla] Girar exactamente un cuarto de vuelta cambia \code{a=} en unas 1024 cuentas. Esto confirma la lectura de 3 bytes del diseño; si los valores no acompañan el giro, anotarlo (sección~\ref{sec:problemas}).
  \item[\casilla] Acercando y alejando el imán, \code{g=} baja y sube; aparecen \code{FUE} muy cerca y \code{DEB} muy lejos.
  \end{itemize}
\end{enumerate}

\section{Paso 2: montaje mecánico y distancia del imán}
\begin{center}
\begin{tikzpicture}[>=Stealth,font=\small]
\fill[gray!25] (-0.6,0) rectangle (0.6,2.2);
\node[left] at (-0.75,1.1) {eje de salida de la reductora};
\fill[red!60] (-0.6,2.2) rectangle (0,2.55);
\fill[blue!55] (0,2.2) rectangle (0.6,2.55);
\node[right] at (0.75,2.38) {imán diametral (N | S)};
\fill[green!40!black] (-1.4,3.35) rectangle (1.4,3.5);
\fill[black] (-0.35,3.2) rectangle (0.35,3.35);
\node[right] at (1.5,3.42) {placa del AS5600 (chip hacia el imán)};
\draw[<->] (-1.0,2.55) -- (-1.0,3.2) node[midway,left] {0,5--3 mm};
\draw[dashed] (0,-0.2) -- (0,3.8) node[above] {eje de giro = centro del chip};
\end{tikzpicture}
\end{center}
\begin{enumerate}
\item El imán va \textbf{pegado en la punta del eje de salida de la caja reductora}, centrado y girando con el eje. Los polos quedan a los costados (imán diametral), no en las caras.
\item El chip queda enfrentado al imán, con el \textbf{centro del chip sobre la prolongación del eje} (desalineación lateral de unas décimas de milímetro como máximo) y las caras paralelas.
\item Distancia imán--chip inicial de 1,5~mm con un separador. Rango admisible 0,5 a 3~mm.
\item La distancia no debe variar al girar: sin juego axial ni excentricidad del eje. Fijar el soporte del sensor a la estructura, no al eje.
\item Cerca del imán, sólo materiales no magnéticos. Evitar tornillos o ejes de acero a pocos milímetros del imán y mantener el sensor lejos del motor 775 y de sus cables.
\item Preferible: elegir para el sensor un eje que gire \textbf{menos de media vuelta a cada lado} del centro en todo el recorrido. Así se puede usar el centro absoluto (paso~4) y no hace falta centrar el yugo al encender.
\end{enumerate}

\textbf{Ajuste fino de la distancia.} Con la Pico encendida y el monitor serie abierto, recorrer el eje de tope a tope:
\begin{itemize}
\item[\casilla] \code{s=OK} en todo el recorrido.
\item[\casilla] \code{g=} cerca de 64 (aceptable aproximadamente entre 30 y 100) y sin grandes variaciones al girar. Si \code{g} sube y baja mucho dentro de una vuelta, el imán está descentrado o el eje tiene excentricidad.
\item[\casilla] Si \code{g} es alto o aparece \code{DEB}, acercar el imán; si es bajo o aparece \code{FUE}, alejarlo.
\end{itemize}

\section{Paso 3: sentido de giro}
\begin{enumerate}
\item Llevar el yugo hacia el lado positivo de la convención del proyecto (roll a la derecha; pitch tirando hacia el piloto).
\item Observar \code{x=} del eje correspondiente: debe crecer hacia positivo.
\item Si crece hacia negativo, cambiar en \code{mis_defines.h} el valor de \code{1} a \code{-1} en
\begin{lstlisting}
ROLL_AS5600_DIRECTION   o   PITCH_AS5600_DIRECTION
\end{lstlisting} Alternativa: pasar el pin DIR de GND a 3,3~V. Usar sólo una de las dos.
\end{enumerate}
El sentido del sensor es independiente del sentido del motor (\code{*_MOTOR_DIRECTION}): no corregir uno con el otro.

\section{Paso 4: centro y extremos}
\subsection{Elegir el tipo de centro de cada eje}
\begin{center}
\begin{tabularx}{\textwidth}{lXX}
\toprule
 & Centro absoluto & Centro al encender\\
\midrule
Condición & El eje del sensor gira como máximo 1800 cuentas (\(\approx 158^\circ\)) hacia cada lado. & Cualquier recorrido, incluso más de una vuelta.\\
Al encender & El yugo puede estar en cualquier posición. & \textbf{El yugo debe estar centrado.}\\
Valor en \code{*_AS5600_CENTER_RAW} & Lectura \code{a=} con el yugo centrado (0..4095). & \code{-1}.\\
Ante una falla larga & Se recupera sola. & Queda en \code{BLQ} hasta reiniciar.\\
\bottomrule
\end{tabularx}
\end{center}
Estimación actual: roll \(\pm 90^\circ\) en el eje del sensor (1024 cuentas, admite centro absoluto); pitch \(\pm 229^\circ\) (2608 cuentas, requiere centro al encender), si el sensor está en el eje de la polea de 75~mm.

\subsection{Medir}
Con el valor de fábrica \code{-1}, encender con el yugo centrado y anotar, para cada eje:
\begin{enumerate}
\item \code{a=} con el yugo centrado: \(a_c\).
\item \code{a=} en el tope negativo: \(a_-\). Llevar el yugo despacio para seguir los cruces de 4095 a 0.
\item \code{a=} en el tope positivo: \(a_+\).
\end{enumerate}
Cuentas a cada extremo, con la diferencia circular (sumar o restar 4096 cuando se cruzó el 0, una vez por cada cruce):
\[
N_+ = \lvert a_+ - a_c \rvert_{\text{circ}}, \qquad N_- = \lvert a_- - a_c \rvert_{\text{circ}}
\]
Ejemplo: \(a_c = 3900\), \(a_+ = 830\) sin vueltas completas: \(N_+ = 830 + 4096 - 3900 = 1026\).

\subsection{Cargar los valores en \texttt{mis\_defines.h}}
\begin{lstlisting}
#define YOKE_POSITION_SOURCE 3        // o compilar con probar.ps1 -As5600
#define ROLL_AS5600_CENTER_RAW  3900  // a_c de roll, o -1
#define PITCH_AS5600_CENTER_RAW -1    // pitch con mas de media vuelta
#define ROLL_AS5600_COUNTS_NEG  1020  // N- medido
#define ROLL_AS5600_COUNTS_POS  1026  // N+ medido
#define PITCH_AS5600_COUNTS_NEG ....  // N- medido
#define PITCH_AS5600_COUNTS_POS ....  // N+ medido
\end{lstlisting}
Si un valor viola un límite (por ejemplo, más de 1800 cuentas con centro absoluto), la compilación se detiene con un mensaje que lo indica.

Compilar con \code{probar.ps1 -As5600}, cargar el UF2 y comprobar:
\begin{itemize}
\item[\casilla] Yugo centrado: \code{x=} cerca de 0.
\item[\casilla] Topes: \code{x=} cerca de \(\pm 10000\) en ambos lados de ambos ejes.
\end{itemize}

\section{Paso 5: ensayo de fallas (motores sin alimentación)}
\begin{enumerate}
\item[\casilla] Desconectar el SDA del sensor de roll: aparece \code{s=I2C} y el comando \code{f=} de \textbf{ambos} ejes pasa a 0.
\item[\casilla] Reconectarlo: con centro absoluto vuelve \code{s=OK} en unos 100~ms; con centro al encender queda \code{s=BLQ} (es lo esperado si estuvo desconectado más de 40~ms) y se recupera reiniciando la Pico con el yugo centrado.
\item[\casilla] Repetir con pitch.
\item[\casilla] Retirar un imán: \code{s=SIN} y fuerza en cero.
\item[\casilla] Arrancar la Pico con un sensor desconectado: la fuerza no se habilita.
\end{enumerate}

\section{Paso 6: tiempo del lazo de control}
La lectura de pitch comparte el bus con el LCD a 100~kHz; el diseño estima pitch 0,6~ms + roll 0,15~ms + LCD 0,3--0,5~ms dentro del ciclo de 2~ms. Debe verificarse:
\begin{enumerate}
\item Con el firmware de potenciómetros (modo 2) y luego con el de AS5600 (modo 3), hacer el mismo vuelo corto con el registro de ensayo (\code{herramientas/ensayos/ENSAYO.cmd}).
\item Comparar en las líneas \code{SC8} el valor \code{lp=} (peor duración del lazo, en \(\mu\)s) y la cantidad de ciclos lentos registrados.
\item[\casilla] Criterio: el modo 3 no aumenta \code{lp=} por encima del umbral de ciclo lento (3000~\(\mu\)s) ni la cantidad de ciclos lentos frente al modo 2.
\end{enumerate}
Si no se cumple, no usar el modo 3 en vuelo e informar los registros para evaluar la lectura asíncrona o separar el LCD a otro ciclo.

\section{Paso 7: primera prueba con motores}
\begin{enumerate}
\item Ganancia total baja desde el menú y el yugo sujeto con la mano.
\item Encender la fuente de los motores con el yugo centrado (siempre, si algún eje usa centro al encender).
\item Verificar que el centrado empuja hacia el centro en ambos sentidos de ambos ejes. Si empuja hacia afuera, revisar primero el sentido del sensor (paso~3) y recién después el del motor.
\item Subir la ganancia gradualmente y comprobar que damper y fricción no generan vibración en reposo (ruido del sensor). Si aparece, informar el valor de \code{g=} y una captura de la telemetría.
\end{enumerate}

\section{Solución de problemas}\label{sec:problemas}
\begin{center}
\begin{tabularx}{\textwidth}{lX}
\toprule
Síntoma & Causa probable y acción\\
\midrule
\code{s=I2C} permanente & Cable SDA/SCL invertido o suelto, módulo sin 3,3~V, GND no común. Revisar continuidad.\\
\code{s=SIN} con el imán bien puesto & Imán no diametral (axial) o módulo que no informa MD. Probar otro imán; si persiste y \code{a=} sí varía con el giro, compilar con \code{YOKE_AS5600_REQUIRE_MD 0}.\\
\code{a=} salta o cambia de sentido al azar & Pin DIR suelto: conectarlo a GND.\\
\code{a=} no acompaña el giro & Anotar los valores observados: puede requerir leer STATUS y ángulo en transacciones separadas (cambio de firmware).\\
\code{g=} varía mucho en una vuelta & Imán descentrado o eje excéntrico. Recentrar.\\
\code{x=} no llega a \(\pm 10000\) o se satura antes del tope & Cuentas a los extremos mal medidas: repetir el paso~4.\\
\code{BLQ} sin desconectar nada & Arranque sin el yugo centrado, sobrerrecorrido mayor que 1,25 veces las cuentas a un extremo, o pausa del lazo mayor de 40~ms. Revisar las cuentas y el tiempo de lazo.\\
LCD falla desde que se conectó pitch & Bus I2C0 sobrecargado o con 5~V: revisar pull-ups y adaptador de nivel.\\
\bottomrule
\end{tabularx}
\end{center}

\section{Registro de resultados}
\begin{center}
\renewcommand{\arraystretch}{2.0}
\begin{tabularx}{\textwidth}{lXX}
\toprule
Dato & Roll & Pitch\\
\midrule
Distancia imán--chip (mm) & & \\
\code{g=} en centro / extremos & & \\
Tipo de centro (absoluto / al encender) & & \\
\(a_c\), \(a_-\), \(a_+\) & & \\
\(N_-\), \(N_+\) (cuentas) & & \\
Sentido (\code{DIRECTION}) & & \\
Ensayo de fallas (OK / observación) & & \\
\code{lp=} modo 2 / modo 3 (\(\mu\)s) & \multicolumn{2}{l}{}\\
\bottomrule
\end{tabularx}
\end{center}

Con estos resultados, informar a Claude para cerrar la validación y llevar a \code{master} la rama \code{mejora/encoder-as5600}.
\end{document}
